\section{Homework Week 5}

\begin{exercise}
    Are the following statements correct?
    \begin{enumerate}
        \renewcommand{\labelenumi}{(\theenumi)}
        \item Any submodule of a semisimple module is also semisimple.
        \item Any semisimple Noetherian module is Artinian.
    \end{enumerate}
\end{exercise}

\begin{proof}
    \textbf{(1)} True. Let $V$ be a semisimple module, so
    \[
        V=\bigoplus_{i\in I} S_i
    \]
    where each $S_i$ is simple. Let $W\leq V$ be a submodule. We show that $W$ is a sum of simple submodules, hence semisimple.

    Consider the set $\mathcal{S}$ of all subsets $J\subseteq I$ such that the sum
    \[
        W+\bigoplus_{j\in J}S_j
    \]
    is direct, i.e.
    \[
        W\cap \bigoplus_{j\in J}S_j=0.
    \]
    By Zorn's lemma, there exists a maximal such subset $J_0$. Set
    \[
        U:=W\oplus \bigoplus_{j\in J_0}S_j.
    \]
    For any $i\in I\setminus J_0$, the maximality of $J_0$ forces $U\cap S_i\neq 0$. Since $S_i$ is simple, we have $S_i\subseteq U$. Therefore $S_i\subseteq U$ for every $i\in I$, so $V=U$. Hence
    \[
        V=W\oplus \bigoplus_{j\in J_0}S_j,
    \]
    and
    \[
        W\cong V\Big/\bigoplus_{j\in J_0}S_j \cong \bigoplus_{i\in I\setminus J_0} S_i,
    \]
    which is semisimple.

    \medskip
    \textbf{(2)} True. Let $V$ be a semisimple Noetherian module. Since $V$ is semisimple,
    \[
        V=\bigoplus_{i\in I}S_i
    \]
    Since $V$ is Noetherian, it is finitely generated. Let $v_1,\dots,v_n$ generate $V$.
    For each $k$, because $V=\bigoplus_{i\in I} S_i$ is a direct sum, the element $v_k$ has only finitely many nonzero components, so there exists a finite subset $I_k\subseteq I$ such that
    \[
        v_k\in \bigoplus_{i\in I_k} S_i.
    \]
    Hence
    \[
        V=\langle v_1,\dots,v_n\rangle \subseteq \bigoplus_{i\in I_1\cup\cdots\cup I_n} S_i.
    \]
    Since the reverse inclusion is obvious, we get
    \[
        V=\bigoplus_{i\in I_1\cup\cdots\cup I_n} S_i,
    \]
    and therefore $V$ is a finite direct sum of simple modules. Each $S_i$ is simple, hence both Noetherian and Artinian. A finite direct sum of Artinian modules is Artinian, so $V$ is Artinian.
\end{proof}

\begin{exercise}
    Let \(V\) be an \(R\)-module with a composition series. Prove that the following conditions are equivalent.
    \begin{enumerate}
        \renewcommand{\labelenumi}{(\theenumi)}
        \item \(V\) is semisimple.
        \item Every irreducible submodule of \(V\) is a direct summand of \(V\).
        \item Every maximal submodule of \(V\) is a direct summand of \(V\).
    \end{enumerate}
\end{exercise}

\begin{proof}
    We prove
    \[
        (1)\Longrightarrow(2)\Longrightarrow(3)\Longrightarrow(1).
    \]

    \medskip
    \noindent\textbf{(1) \(\Longrightarrow\) (2).}
    If \(V\) is semisimple, then every submodule of \(V\) is a direct summand.
    In particular, every irreducible submodule of \(V\) is a direct summand.

    \medskip
    \noindent\textbf{(2) \(\Longrightarrow\) (3).}
    Let \(M\) be a maximal submodule of \(V\).
    Then \(V/M\) is simple.

    Since \(M\) is a maximal submodule of \(V\), the quotient \(V/M\) is simple.
    Now let \(S\) be a submodule of \(V\) such that \(S \nsubseteq M\).
    Then \(S+M \neq M\), so
    \[
        (S+M)/M \neq 0.
    \]
    By the second isomorphism theorem,
    \[
        S/(S\cap M)\cong (S+M)/M.
    \]
    Since \((S+M)/M\) is a nonzero submodule of the simple module \(V/M\), it follows that
    \((S+M)/M\) is simple. Hence \(S/(S\cap M)\) is simple.

    Choose \(S\subseteq V\) minimal subject to \(S\nsubseteq M\).
    Then \(S\cap M\) is a maximal submodule of \(S\), and by minimality of \(S\),
    the quotient \(S/(S\cap M)\) is simple.
    But \(S\cap M\) is a proper submodule of \(S\), so by minimality of \(S\), it must be \(0\).
    Thus \(S\) is simple.

    Therefore \(S\) is an irreducible submodule of \(V\) with \(S\nsubseteq M\).
    By (2), \(S\) is a direct summand of \(V\), say
    \[
        V=S\oplus T.
    \]
    Since \(S\nsubseteq M\) and \(S\) is simple, we have \(S\cap M=0\).
    Also \(M\) is maximal and \(S\not\subseteq M\), so
    \[
        M+S=V.
    \]
    Hence
    \[
        V=M\oplus S,
    \]
    and therefore \(M\) is a direct summand of \(V\).

    \medskip
    \noindent\textbf{(3) \(\Longrightarrow\) (1).}
    Let
    \[
        V=V_0\supsetneq V_1\supsetneq \cdots \supsetneq V_n=0
    \]
    be a composition series of \(V\).
    We argue by induction on \(n\), the composition length of \(V\).

    If \(n=1\), then \(V\) is simple, hence semisimple.

    Assume \(n>1\), and suppose the statement holds for modules of composition length \(<n\).
    Since \(V/V_1\) is simple, \(V_1\) is a maximal submodule of \(V\).
    By (3), \(V_1\) is a direct summand of \(V\); thus
    \[
        V=V_1\oplus S
    \]
    for some simple submodule \(S\cong V/V_1\).

    Now \(V_1\) has composition length \(n-1\).
    We claim that every maximal submodule of \(V_1\) is a direct summand of \(V_1\).
    Indeed, let \(N\) be a maximal submodule of \(V_1\).
    Then
    \[
        V/N \cong (V_1/N)\oplus S,
    \]
    and since \(V_1/N\) is simple, \(N\oplus S\) is a maximal submodule of \(V\).
    By (3), \(N\oplus S\) is a direct summand of \(V\), say
    \[
        V=(N\oplus S)\oplus T.
    \]
    Because
    \[
        V=V_1\oplus S,
    \]
    intersecting with \(V_1\) gives
    \[
        V_1=N\oplus T.
    \]
    So \(N\) is a direct summand of \(V_1\).

    Hence \(V_1\) satisfies condition (3), and by the induction hypothesis \(V_1\) is semisimple.
    Therefore
    \[
        V=V_1\oplus S
    \]
    is semisimple as well.

    This proves \((3)\Longrightarrow(1)\), and therefore the three conditions are equivalent.
\end{proof}

\begin{exercise}
    Prove that
    \[
        \operatorname{Soc}(U \oplus V)=\operatorname{Soc}(U)\oplus \operatorname{Soc}(V).
    \]
\end{exercise}

\begin{proof}
    We first show
    \[
        \operatorname{Soc}(U)\oplus \operatorname{Soc}(V)\subseteq \operatorname{Soc}(U\oplus V).
    \]
    Indeed, \(\operatorname{Soc}(U)\) and \(\operatorname{Soc}(V)\) are semisimple, so
    \(\operatorname{Soc}(U)\oplus \operatorname{Soc}(V)\) is a semisimple submodule of
    \(U\oplus V\). Hence it is contained in \(\operatorname{Soc}(U\oplus V)\).

    For the reverse inclusion, let \(S\) be a simple submodule of \(U\oplus V\), and let
    \[
        \pi_U:U\oplus V\to U, \qquad \pi_V:U\oplus V\to V
    \]
    be the canonical projections. Consider the restrictions
    \[
        \pi_U|_S:S\to U, \qquad \pi_V|_S:S\to V.
    \]
    Since \(S\) is simple, the image of any homomorphism defined on \(S\) is either \(0\) or
    simple. Therefore \(\pi_U(S)\) is either \(0\) or a simple submodule of \(U\), and similarly
    \(\pi_V(S)\) is either \(0\) or a simple submodule of \(V\). Hence
    \[
        \pi_U(S)\subseteq \operatorname{Soc}(U),
        \qquad
        \pi_V(S)\subseteq \operatorname{Soc}(V).
    \]

    Now for any \(s\in S\), we can write
    \[
        s=(\pi_U(s),\pi_V(s)).
    \]
    Since \(\pi_U(s)\in \operatorname{Soc}(U)\) and
    \(\pi_V(s)\in \operatorname{Soc}(V)\), it follows that
    \[
        s\in \operatorname{Soc}(U)\oplus \operatorname{Soc}(V).
    \]
    Thus
    \[
        S\subseteq \operatorname{Soc}(U)\oplus \operatorname{Soc}(V).
    \]

    Since \(\operatorname{Soc}(U\oplus V)\) is the sum of all simple submodules of \(U\oplus V\),
    we obtain
    \[
        \operatorname{Soc}(U\oplus V)\subseteq \operatorname{Soc}(U)\oplus \operatorname{Soc}(V).
    \]

    Combining the two inclusions, we conclude that
    \[
        \operatorname{Soc}(U\oplus V)=\operatorname{Soc}(U)\oplus \operatorname{Soc}(V).
    \]
\end{proof}

\begin{exercise}
    Let \(R\) be an arbitrary ring. Then every two-sided ideal of the matrix ring \(M_n(R)\) can be expressed in the form \(M_n(I)\), where \(I\) is a two-sided ideal of \(R\). In particular, \(R\) is simple if and only if \(M_n(R)\) is simple.

    \textit{Hint.} For an ideal \(L\) of \(M_n(R)\), let \(I\) be the set of \((1,1)\)-entries of \(L\). Show that \(I\) is an ideal of \(R\) and \(L=M_n(I)\).
\end{exercise}

\begin{proof}
    Let $E_{ij}$ denote the matrix with $1$ in position $(i,j)$ and $0$ elsewhere. For a two-sided ideal $L$ of $M_n(R)$, define
    \[
        I:=\{a\in R\mid aE_{11}\in L\}.
    \]
    We first note that for any $A=(a_{kl})\in L$, we have
    \[
        E_{1k}\cdot A\cdot E_{l1}=a_{kl}\,E_{11}\in L,
    \]
    so every entry of every matrix in $L$ belongs to $I$. In particular,
    \[
        I=\{a_{kl}\mid A=(a_{kl})\in L,\; 1\leq k,l\leq n\}.
    \]

    \textbf{$I$ is a two-sided ideal of $R$.}
    If $a,b\in I$, then $aE_{11},bE_{11}\in L$, so $(a-b)E_{11}=aE_{11}-bE_{11}\in L$, hence $a-b\in I$. For $r\in R$, we have $(rE_{11})\cdot(aE_{11})=raE_{11}\in L$ and $(aE_{11})\cdot(rE_{11})=arE_{11}\in L$, so $ra,ar\in I$. Thus $I$ is a two-sided ideal of $R$.

    \textbf{$L=M_n(I)$.}
    Since every entry of every matrix in $L$ lies in $I$, we have $L\subseteq M_n(I)$. Conversely, let $aE_{ij}\in M_n(I)$ with $a\in I$. Then $aE_{11}\in L$, and
    \[
        aE_{ij}=E_{i1}\cdot(aE_{11})\cdot E_{1j}\in L.
    \]
    Since every matrix in $M_n(I)$ is a sum of such elements $aE_{ij}$, we get $M_n(I)\subseteq L$. Therefore $L=M_n(I)$.

    \textbf{$R$ is simple $\Longleftrightarrow M_n(R)$ is simple.}
    The map $I\mapsto M_n(I)$ is a bijection from the set of two-sided ideals of $R$ to the set of two-sided ideals of $M_n(R)$. Hence $R$ has no nontrivial two-sided ideals if and only if $M_n(R)$ has no nontrivial two-sided ideals.
\end{proof}

\begin{exercise}
    Let \(W\) be a submodule of a semisimple module \(V\). Show that \(W\) and \(V/W\) are also semisimple.
\end{exercise}

\begin{proof}
    Since $V$ is semisimple, write $V=\bigoplus_{i\in I}S_i$ with each $S_i$ simple.

    \textbf{$W$ is semisimple.}
    By Zorn's lemma, choose a maximal subset $J\subseteq I$ such that $W\cap\bigoplus_{j\in J}S_j=0$. Set $U=W\oplus\bigoplus_{j\in J}S_j$. For any $i\in I\setminus J$, the maximality of $J$ gives $U\cap S_i\neq 0$, and since $S_i$ is simple, $S_i\subseteq U$. Therefore $V\subseteq U$, so $V=U=W\oplus\bigoplus_{j\in J}S_j$. Thus
    \[
        W\cong V\Big/\bigoplus_{j\in J}S_j\cong \bigoplus_{i\in I\setminus J}S_i,
    \]
    which is a direct sum of simple modules, hence semisimple.

    \textbf{$V/W$ is semisimple.}
    From the decomposition $V=W\oplus\bigoplus_{j\in J}S_j$, we obtain
    \[
        V/W\cong \bigoplus_{j\in J}S_j,
    \]
    which is a direct sum of simple modules, hence semisimple.
\end{proof}

